\newpage
\thispagestyle{empty}
\chapter*{\sffamily Resumen}
\addcontentsline{toc}{chapter}{Resumen}%
\begin{center}
\textbf{\large \thesisname}
\end{center}
\par Este trabajo presenta el diseño de un kit educativo para la clasificación de objetos mediante visión por computador, empleando el robot manipulador PhantomX Pincher, la plataforma Raspberry Pi 5 y el sistema operativo robótico ROS2. El proyecto incluyó la identificación de requerimientos de hardware y software, el diseño conceptual de la estructura física del kit con componentes fabricados mediante impresión 3D, y la propuesta de una arquitectura de software modular basada en ROS2. Se desarrollaron simulaciones funcionales en Gazebo y se implementó un workspace completo que integra MoveIt2 para planificación de movimiento y OpenCV para procesamiento de imágenes. Se fabricaron siete prototipos físicos del kit y se generaron tres repositorios digitales con imagen del sistema operativo, modelado 3D completo y workspace de software con documentación educativa. El kit desarrollado representa un ahorro significativo respecto a alternativas comerciales. La validación con estudiantes confirmó la viabilidad del sistema como herramienta pedagógica para enseñanza de robótica, visión por computador y ROS2.
\\[2cm]
\textbf{Palabras clave:} \palabrasclave

\newpage
\thispagestyle{empty}
\chapter*{\sffamily Abstract}
\addcontentsline{toc}{chapter}{Abstract}%
\begin{center}
\textbf{\large \thesisnameeng}
\end{center}
\par This work presents the design of an educational kit for object classification using computer vision, employing the PhantomX Pincher robotic manipulator, the Raspberry Pi 5 platform, and the ROS2 robotic operating system. The project included the identification of hardware and software requirements, the conceptual design of the kit's physical structure with components manufactured through 3D printing, and the proposal of a modular software architecture based on ROS2. Functional simulations were developed in Gazebo, and a complete workspace integrating MoveIt2 for motion planning and OpenCV for image processing was implemented. Seven physical prototypes of the kit were manufactured, and three digital repositories were generated with the operating system image, complete 3D modeling, and software workspace with educational documentation. The developed kit represents significant savings compared to commercial alternatives. Validation with students confirmed the system's viability as a pedagogical tool for teaching robotics, computer vision, and ROS2.
\\[2cm]
\textbf{Keywords:} \keywords

%Usar si es requerido, de no necesitarse favor comentar con % las lineas
\newpage 
\thispagestyle{empty}
\chapter*{\sffamily Zusammenfassung}
\addcontentsline{toc}{chapter}{Zusammenfassung}%
\begin{center}
\textbf{\large \thesisnamelang}
\end{center}
\par Diese Arbeit präsentiert das Design eines Bildungs-Kits zur Objektklassifizierung mittels Computersehen unter Verwendung des PhantomX Pincher Robotermanipulators, der Raspberry Pi 5 Plattform und des ROS2 Roboter-Betriebssystems. Das Projekt umfasste die Identifizierung von Hardware- und Software-Anforderungen, das konzeptionelle Design der physischen Struktur des Kits mit Komponenten aus 3D-Druck, sowie den Vorschlag einer modularen Software-Architektur basierend auf ROS2. Funktionale Simulationen wurden in Gazebo entwickelt und ein vollständiger Workspace implementiert, der MoveIt2 für Bewegungsplanung und OpenCV für Bildverarbeitung integriert. Sieben physische Prototypen des Kits wurden hergestellt und drei digitale Repositories generiert mit Betriebssystem-Image, vollständiger 3D-Modellierung und Software-Workspace mit pädagogischer Dokumentation. Das entwickelte Kit stellt eine signifikante Ersparnis im Vergleich zu kommerziellen Alternativen dar. Die Validierung mit Studenten bestätigte die Eignung des Systems als pädagogisches Werkzeug für den Unterricht in Robotik, Computersehen und ROS2.
\\[2cm]
\textbf{Schlüsselwörter:} \schlusselworter

